% APPROVAL ONLY — NOT A COMBINED PATCH. Alternatives A and B are mutually exclusive.
% Actual latest RevThree/RevFive wrappers were not recovered. Keep existing wrappers unchanged.
% Fragments are deliberately unwrapped. Do not input this whole file into a manuscript.

% Fragment: section3_weighted_medoid — A: describe traced historical implementation
% Source lines 501–509 in the designated 12 October snapshot.
\[
m_i
\in
\arg\min_{u\in U_i}
\sum_{\ell\in U_i}\pi_\ell
\|\bm c_u-\bm c_\ell\|,
\qquad
\bm r_i=\bm c_{m_i}.
\]

% Fragment: section3_assignment — A: describe traced historical implementation
% Source lines 545–556 in the designated 12 October snapshot.
 The sample units are then matched to the medoid reference points by solving the
assignment problem
\[
\tau
=
\arg\min_{\tau}
\sum_{i=1}^n
\|\bm c_{\tau(i)}-\bm r_i\|,
\qquad
\tau(i)\in s,
\]
where \(\tau\) is a one-to-one mapping from clusters to sampled units. Each

% Fragment: section3_historical_kernel — A: describe traced historical implementation
% Source lines 586–597 in the designated 12 October snapshot.
where, for the traced historical computations,
\(\mathbf H=h^2\mathbf I_d\) with \(h=N^{-1/(d+4)}\), the isotropic
Scott bandwidth. The spherical uniform (\texttt{tophat}) kernel is
\[
K(\bm u)=\frac{1}{v_d}\mathbf 1\{\|\bm u\|_2<1\},
\qquad v_d=\frac{\pi^{d/2}}{\Gamma(1+d/2)}.
\]
The same numerical bandwidth is used for the original and translated
coordinates, without covariance scaling or renormalizing the translated
population. The historical KDE evaluation used relative tolerance
\(10^{-4}\). Applying the same estimator to the translated coordinates

% Fragment: section4_replications — Both A and B
% Source lines 977–978 in the designated 12 October snapshot.
\citep{sampling_pkg}. The reported summaries combine simulation stages.
The auxiliary-index Monte Carlo summaries were inherited from earlier
output files recording \(5{,}000\) replications per setting. The subsequent
medoid-DI update requested \(2{,}000\) samples per setting. Retained outputs
label \(106\) setting summaries with this count; \(104\) reproduce the
reported DI mean and standard deviation, whereas two reported means differ.
For \(32\) other Monte Carlo settings, the count underlying the final DI
summary cannot be established from retained execution records. The nominal
counts are supported by source and summary metadata rather than independently
retained per-replication completion logs. The NMS/GMS DI comparisons use
probability-weighted finite-design enumeration, with the number of enumerated
entries reported separately for each design; they are not Monte Carlo
replications. Accordingly, a common \(5{,}000\)-replication count for every
reported statistic is not supported by the retained records.

% Fragment: algorithm2_medoid — A: describe traced historical implementation
% Source lines 1490–1505 in the designated 12 October snapshot.
\State  For each cluster \(U_i\), define the reference point \(\bm q_i\) by
\[
\bm q_i
=
\begin{cases}
\bm c_{m_i}, &
\rho=\textsc{medoid},\quad
m_i\in\displaystyle\arg\min_{u\in U_i}
\sum_{\ell\in U_i}\pi_\ell\|\bm c_u-\bm c_\ell\|, \\[2ex]
\displaystyle
\frac{\sum_{\ell\in U_i}\pi_\ell\bm c_\ell}
{\sum_{\ell\in U_i}\pi_\ell}, &
\rho=\textsc{centroid}.
\end{cases}
\]
In the traced historical medoid-update implementation, the first minimizing unit in cluster order was selected.

% Fragment: algorithm2_assignment — A: describe traced historical implementation
% Source lines 1508–1518 in the designated 12 October snapshot.
\State Match the reference points to the sampled units by solving
\[
\tau
=
\arg\min_{\tau}
\sum_{i=1}^n
\|\bm c_{\tau(i)}-\bm q_i\|,
\qquad
\tau(i)\in s,
\]
where \(\tau\) is one-to-one.

% Fragment: comment5_translation_qualification — Both A and B
% Source lines 901–917 in the designated 12 October snapshot.
    \item[(iii)] Conditional on a fixed hard partition, fixed reference
points, and a fixed medoid-to-sample pairing, DI is invariant to a common
translation of the matched sample coordinates. Indeed, replacing every
matched sample point by \(\bm c_{\tau(i)}+\bm a\) replaces every transformed
population point by \(\widetilde{\bm c}_\ell+\bm a\). A fixed-bandwidth,
translation-invariant KDE has the same value at each corresponding
transformed unit, leaving the DI formula unchanged. This result does not
guarantee invariance after sample-adaptive reclustering or rematching.
Ordinary Euclidean optimal matching can change when only the sample is
translated. Thus, for that historical objective, unchanged matching is an
additional condition, not a general consequence of unchanged reference points.
The numerical translation examples establish the behavior of the saved
configurations only. DI is not generally rotation invariant; the illustrated
rotated targets are projected to distinct population-grid units.

% Fragment: section3_keep_squared_alternative — B: retain squared definition; explicitly separate historical pipeline
% Source lines 545–556 in the designated 12 October snapshot.
The formal definition below retains squared-distance representatives,
squared Euclidean matching, and the box kernel specified here. The traced
historical Monte Carlo computations instead used weighted preliminary
\(k\)-means followed by probability refinement, ordinary-distance medoids,
ordinary Euclidean matching, and a spherical \texttt{tophat} KDE.
The retained results must therefore be identified as results of that
historical variant, rather than numerical validation of the present
OT/squared-distance/box-kernel specification. Runtime provenance remains
unresolved for the other blocks identified in the reproducibility audit.

 The sample units are then matched to the medoid reference points by solving the
assignment problem
\[
\tau
=
\arg\min_{\tau}
\sum_{i=1}^n
\|\bm c_{\tau(i)}-\bm r_i\|^2,
\qquad
\tau(i)\in s,
\]
where \(\tau\) is a one-to-one mapping from clusters to sampled units. Each

% Fragment: historical_clustering_disclosure — Both A and B; essential historical-method disclosure
% Source lines 488–489 in the designated 12 October snapshot.
\(\{\bm c_\ell:\ell\in s\}\) as initial centers, the formal definition applies
the hard IP\(n\)M procedure to obtain \(n\) spatially compact clusters.
The traced historical Monte Carlo implementation used weighted preliminary
\(k\)-means followed by open-TSP ordering and inclusion-probability refinement,
rather than the entropic OT preliminary stage described in Section~2.
Its retained results are distinguished from computations of the formal OT
variant; exact runtime attribution is unresolved for the other audited blocks.

% Fragment: abstract_translation — Both A and B
% Source lines 143–143 in the designated 12 October snapshot.
bounded measure with a conditional translation property that distinguishes under-dispersion,

% Fragment: introduction_translation — Both A and B
% Source lines 233–235 in the designated 12 October snapshot.
interpreted as approximately well spread. With the reference partition and matching held fixed, a common translation of the matched sample coordinates leaves DI unchanged. This conditional property concerns the relative arrangement of the matched sample and does not guarantee invariance after sample-adaptive reclustering or rematching.

% Fragment: conclusion_translation — Both A and B
% Source lines 1120–1121 in the designated 12 October snapshot.
The proposed index is bounded and directly interpretable, and has the conditional translation property described in Section~\ref{Sec:Measure}. Negative values indicate under-dispersion or local clustering,

% Fragment: appendix_translation_example — Both A and B
% Source lines 1150–1151 in the designated 12 October snapshot.
In panel~(b), a uniform translation is applied to all sample units and DI
remains unchanged for this example. This observation alone does not establish invariance after general sample-adaptive reclustering or rematching.

% Fragment: section3_historical_medoid_ties — A: describe the historical notebook tie convention
% Source lines 510–512 in the designated 12 October snapshot.
In the traced historical medoid-update implementation, the first minimizing
unit in cluster order was selected.

% Fragment: algorithm2_historical_preliminary_stage — A only: describe the historical pipeline; do not mix with the formal OT version
% Source lines 1485–1487 in the designated 12 October snapshot.
\State For the historical DI variant, initialize weighted \(k\)-means
at the observed sample units and apply the historical open-TSP ordering and
inclusion-probability refinement to construct the reference partition.
\label{line:upbkm}

% Optional Comment5 explanatory paragraph. Preserve the actual RevFive wrapper.
For a fixed reference and a full, unweighted one-to-one assignment,
let \(Q_{\bm a}(\tau)=\sum_{i=1}^n
\|\bm c_{\tau(i)}+\bm a-\bm r_i\|^2\). Then
\[
Q_{\bm a}(\tau)=Q_{\bm0}(\tau)
+2\bm a^\top\left(\sum_{\ell\in s}\bm c_\ell-
\sum_{i=1}^n\bm r_i\right)+n\|\bm a\|^2.
\]
The added term is independent of \(\tau\), so the minimizing permutation
set is unchanged. A consistent tie rule is needed to preserve the selected
permutation. This expansion applies to squared distances and a full
bijection; it does not apply in general to ordinary Euclidean matching or
to rectangular matching that may select a different sample subset.
With ordinary distances, fixed reference points alone do not guarantee
unchanged optimal matching. In either case, the conditional DI result also
requires an unchanged partition and a fixed bandwidth evaluated at the
corresponding transformed units.
